\documentclass[13pt, handout]{beamer}

\usetheme{Madrid}
\usecolortheme{beaver}
\setbeamertemplate{navigation symbols}{}
\setbeamertemplate{page number in head/foot}{}
\setbeamercolor*{item}{fg=red}

\usepackage[portuguese]{babel}
\usepackage{tikz}
\usepackage{ragged2e}

\title{Gestão de uma frota de táxis — Algoritmos de procura para otimização}
\subtitle{{}\color{black}Inteligência Artificial}
\author[Grupo 33]{
    \begin{tabular}{lc}
         André Santos & A106854 \\
         Daniel Parente & A107363 \\
         Pedro Ferreira & A107292 \\
         Tiago Martins & A106938
    \end{tabular}
}

\institute[DIUM]{
    Universidade do Minho -- Licenciatura em Engenharia Informática
}
\date[2026/01/15]{\scriptsize 15 de janeiro de 2026}

\begin{document}
\frame{\titlepage}

\begin{frame}{Índice}
    \tableofcontents
\end{frame}

\justifying

\section{Descrição do Problema}

\begin{frame}{Descrição do Problema -- Enquadramento e Objetivo }

    Tem-se:
    \begin{itemize}
        \item Problema de decisão em tempo real, com restrições
        \item Espaço de dados dinâmico
    \end{itemize}

    \vspace{0.5cm}

    Para isso:
    \begin{itemize}
        \item Formulou-se como problema de procura
        \item Implementaram-se vários algoritmos
        \item Analisaram-se os resultados
    \end{itemize}
    
\end{frame}

\section{Formulação do Problema}

\begin{frame}{Formulação do Problema -- Tipo de Problema}

    \begin{itemize}
        \item Procura em grafos
        \item Não-determinístico
        \item Ambiente dinâmico
        \item Multi-agente
        \item Parcialmente observável
    \end{itemize}
    
\end{frame}

\begin{frame}{Formulação do Problema -- Estado do Sistema}

    \[
        E = \langle T, V, P, B \rangle
    \]

    \begin{itemize}
        \item T: tempo acumulado
        \item V: estado dos veículos
        \item P: pedidos pendentes
        \item B: pontos de abastecimento
    \end{itemize}
    
\end{frame}

\begin{frame}{Formulação do Problema -- Estado Inicial e Objetivo}

    Estado Inicial:
    \begin{itemize}
        \item Frota disponível
        \item Autonomia máxima
        \item Sem pedidos inicias
    \end{itemize}

    \vspace{0.5cm}

    Objetivo:
    \begin{itemize}
        \item Sem pedidos pendentes
        \item Sem viagens ativas
    \end{itemize}
    
    \end{frame}

\begin{frame}{Formulação do Problema -- Operadores}

    \begin{itemize}
        \item Atribuir pedido
        \item Mover veículo
        \item Transportar passageiros
        \item Abastecer/recarregar
        \item Esperar
    \end{itemize}
    
\end{frame}

\begin{frame}{Formulação do Problema -- Função de Custo}

    \[
        C = w_t \times T + w_m \times M + w_e \times E + w_p \times P
    \]

    \begin{itemize}
        \item T: Tempo
        \item M: Custo monetário
        \item E: Emissões
        \item P: Penalizações
    \end{itemize}
\end{frame}

\section{Modelação e Arquitetura}

\begin{frame}{Modelação e Arquitetura -- Mapa}

    \begin{itemize}
        \item Uso de OSMnx
        \item Grafo ponderado de várias cidades de Braga, Portugal
    \end{itemize}
    \vspace{0.5cm}

    Nós: Ruas, estações, postos de abastecimento, ...
    \vspace{0.5cm}
    
    Arestas: distância, tempo, velocidade
    \vspace{0.5cm}
    
    Zonamento: centro vs periferia, zonas de trânsito intenso
    
\end{frame}

\begin{frame}{Modelação e Arquitetura -- Sistema}

    Separação clara:
    \begin{itemize}
        \item Simulador
        \item Supervisor
        \item Algoritmos
        \item Estado/Estruturas de Dados
        \item Vistas
    \end{itemize}

    \vspace{0.5cm}

    Permite:
    \begin{itemize}
        \item Isolamento da lógica de procura
        \item Simulação controlada
        \item Expansibilidade futura
    \end{itemize}
    
\end{frame}

\section{Algoritmos de Procura}

\begin{frame}{Algoritmos de Procura -- Não Informados}

        DFS:
        \begin{itemize}
            \item sem garantias
            \item apenas para comparações académicas
        \end{itemize}

        \vspace{0.5cm}
        
        BFS:
        \begin{itemize}
            \item minimiza saltos, mas...
            \item ignora pesos
            \item inviável em situações reais
        \end{itemize}
\end{frame}

\begin{frame}{Algoritmos de Procura -- Informados}

        Heurística: Distância Euclidiana
        \vspace{0.5cm}
        
        Greedy:
        \begin{itemize}
            \item extremamente rápido
            \item não ótimo, mas...
            \item útil quando a velocidade de decisão é crítica
        \end{itemize}

        \vspace{0.5cm}
        A* (A-Star):
        \begin{itemize}
            \item núcleo do sistema
            \item soma ponderada de três fatores: tempo, custo e emissões
            \item adaptável a diferentes objetivos e contextos
        \end{itemize}
\end{frame}

\section{Simulação dinâmica e métricas}

\begin{frame}{Simulação dinâmica e métricas}

    Dinâmica simulada:
    \begin{itemize}
        \item pedidos aleatórios ao longo do tempo
        \item trânsito variável
        \item autonomia decrescente
        \item filas em postos
        \item falhas pontuais
    \end{itemize}

    \vspace{0.5cm}
    Métricas recolhidas:
    \begin{itemize}
        \item tempo médio de resposta
        \item ocupação da frota
        \item custo total
        \item pedidos não atendidos
    \end{itemize}
    
\end{frame}

\section{Conclusões e Demonstração do Sistema}

\begin{frame}{Conclusões e Demonstração do Sistema}
    
    Apesar de algumas dificuldades, temos:
    \begin{itemize}
        \item Problema bem definido
        \item Modelação apropriada e realista
        \item Algoritmos clássicos bem desenvolvidos
        \item Sistema expansível
    \end{itemize}

    \vspace{0.5cm}
    Demonstração do Sistema!
    
\end{frame}

\frame{\titlepage}

\end{document}
